\section{Methods}
\label{sec:methods}

\subsection{Study design}

We evaluate whether a selective prediction policy calibrated at one hospital
remains reliable at another when pre-diagnostic clinical context changes in
availability or alignment. The design is a retrospective pre-deployment
evaluation.

Two shifts are crossed:

\begin{itemize}
  \item \textbf{Institutional shift} --- source site versus external site.
  \item \textbf{Context shift} --- three controlled interventions on
    pre-diagnostic text.
\end{itemize}

Crossing them is the point. Varying one while holding the other fixed cannot
recover their interaction, and the interaction is what determines whether a
model's measured benefit at one site is also its exposure at the next.

\subsection{Sites and roles}

\begin{table}[htbp]
\centering
\caption{Site roles. The external site is locked: architecture selection, early
stopping, temperature selection, threshold selection, text-cleaning rules and
context-state rules may not be revised using it. These prohibitions are encoded
in the protocol registry and enforced by an automated validator.}
\label{tab:sites}
\begin{tabular}{llll}
\toprule
Role & Dataset & Function & Tuning \\
\midrule
Source   & MIMIC-CXR-JPG v2.1.0~\cite{johnson2019mimiccxr}
         & development and internal evaluation & permitted \\
External & CheXpert Plus~\cite{chexpertplus2024}
         & locked external evaluation only & prohibited \\
\bottomrule
\end{tabular}
\end{table}

\subsection{Data intake and integrity}

Source metadata and the report archive were verified against the official
SHA-256 manifest before use. The report archive was checked entry by entry:
\num{227835} members, zero CRC mismatches, zero read errors, zero undecodable
members. Three independent sources agree exactly on the study inventory --- the
study list, the report archive members, and the checksum manifest --- and the
record list agrees exactly with the manifest's DICOM entries (\num{377110}
records). Patient disjointness across the official splits holds exactly.

All \num{193282} source frontal images were acquired and hashed against the
publisher manifest: \num{193282} of \num{193282} matched. One file failed on
first pass as a partial object left by an interrupted transfer and was
re-fetched and re-verified.

At the external site, \num{191071} frontal images were acquired and checked
against the publisher's per-file checksums, with zero mismatches. One image
matches its recorded checksum yet does not decode, so the damage is upstream;
it is excluded, leaving \num{191070} usable images
(\cref{sec:limitations}).

Verifying that data \emph{arrived} is not the same as verifying that it
\emph{loads}. The checksum pass accepted the corrupt external file, and only a
full decode revealed it.

\subsection{Report structure and the permitted text policy}
\label{sec:report-structure}

The protocol requires model input to be \emph{pre-diagnostic}: text available
before the interpretation exists. Findings, impression, and full reports are
prohibited as inputs and may act only as label sources.

All \num{227835} source reports were parsed structurally using a curated
section-header vocabulary and header-driven --- not position-driven ---
extraction. Permitted context is \texttt{indication} plus \texttt{history}, the
structural analogue of the external site's two permitted sections. Widening the
set was measured and rejected: adding \texttt{examination} gains
\SI{0.11}{\percent} of coverage, while adding \texttt{comparison} gains
\SI{1.49}{\percent} but admits text that can quote prior diagnostic conclusions.

Three structural hazards were identified and are handled explicitly.

\begin{enumerate}
  \item \textbf{Preliminary interpretations embedded in reports.}
    \num{17551} reports (\SI{7.70}{\percent}) contain a \texttt{WET READ}
    section --- a preliminary radiologist reading. This is post-diagnostic text
    sitting inside the report body. Any extraction rule that takes ``everything
    before FINDINGS'' as context silently ingests it. We exclude these sections
    while retaining the study.
  \item \textbf{Non-monotone section order.} \num{181} reports place a
    pre-diagnostic section \emph{after} the first diagnostic section, so
    position-based splitting mis-assigns their text. Extraction is therefore
    header-driven.
  \item \textbf{Non-separable label sections.} \num{134} reports merge findings
    and impression into one section, and \num{66} expose no recognised header.
    Both are excluded (\num{200} studies, \SI{0.09}{\percent}).
\end{enumerate}

The first hazard is the most transferable: it applies to any pipeline that
treats MIMIC report text as pre-diagnostic context, not only to this one.

\subsection{Context states and interventions}

Natural states are assigned by an informativeness rule defined on the source
site. An \emph{effective token} is a whitespace-delimited token containing at
least one alphanumeric character, counted after de-identification placeholder
runs are removed, so a body of placeholders does not count as context.

\begin{table}[htbp]
\centering
\caption{Natural context states. $T=3$ is primary; $T=5$ and $T=10$ are
pre-specified sensitivity analyses. The rule is fixed on the source site and
transferred to the external site without refitting.}
\label{tab:states}
\begin{tabular}{lll}
\toprule
State & Name & Rule \\
\midrule
N1 & informative      & effective tokens $\geq T$ \\
N2 & low information  & $0 <$ effective tokens $< T$ \\
N3 & absent           & effective tokens $= 0$ \\
\bottomrule
\end{tabular}
\end{table}

Choosing $T$ is consequential and was therefore pre-specified. Across the swept
range the confirmatory cohort spans \num{187847} studies at $T=1$ to \num{8996}
at $T=30$ --- a twentyfold range driven by a parameter no prior specification
fixes. We report the full sensitivity curve.

Three controlled interventions are applied to N1 studies only, since removing
context that was never present, or misaligning context that was already absent,
tests nothing:

\begin{itemize}
  \item \textbf{C0, original context} --- identity.
  \item \textbf{C1, no context} --- all permitted context replaced by
    \texttt{[NO\_CONTEXT]}.
  \item \textbf{C2, misaligned context} --- deterministic permutation under
    constraints: same institution, same split, different patient, same
    context-length bin, seed \num{20260718}.
\end{itemize}

C2 is implemented as a rotation within strata rather than a shuffle. A shuffle
can leave fixed points, and a fixed point is silently a C0 study sitting inside
C2, which would bias the misalignment effect toward zero without any visible
failure. Five invariants are asserted at run time rather than assumed, including
that every C2 recipient holds a different patient's text and that image-only
predictions are identical across all three conditions.

\subsection{Cohort construction and evaluation partition}

\begin{table}[htbp]
\centering
\caption{Source-site cohort construction.}
\label{tab:cohort}
\begin{tabular}{lrr}
\toprule
Step & Removed & Remaining \\
\midrule
All studies                      & ---          & \num{227835} \\
No recognised section            & \num{66}     & \num{227769} \\
Merged label sections            & \num{133}    & \num{227636} \\
Frontal view required            & \num{9688}   & \num{217948} \\
Official split assigned          & \num{0}      & \num{217948} \\
Context informative ($T=3$)      & \num{15001}  & \num{202947} \\
Impression label source present  & \num{29375}  & \textbf{\num{173572}} \\
\bottomrule
\end{tabular}
\end{table}

The final source cohort is \num{173572} studies, \num{58689} patients,
\num{193282} frontal images. The official test split yields only \num{282}
patients within this cohort, which a design-stage power analysis showed is
materially underpowered for the primary hypothesis, so a prespecified
patient-disjoint evaluation partition was carved from the training split.
Assignment used patient identifiers and the frozen seed only --- no label,
outcome, image, or model output participated --- so this is a design
construction, not a data-dependent selection. All tiers are verified
patient-disjoint, and reserved tiers are unreadable in code by any stage that
could influence a model, a threshold, or a policy.

\subsection{Models and selective policy}

Four models: image-only (M1, non-multimodal control), text-only (M2,
text-signal and shortcut control), probability late fusion (M3), and
feature-concatenation fusion (M4). M3 holds no trainable parameters by
construction, which is what allows any advantage M4 shows over it to be
attributed to the learned fusion rather than to the mere presence of two
modalities.

Backbones are DenseNet-121 initialised from ImageNet at \SI{224}{px} for the
image arm and BERT-base-uncased at 128 tokens for the text arm, fixed by
amendment before any model was trained. The choice is deliberately conventional:
the study claims no architectural novelty, so the backbone is a control held
constant rather than a variable to optimise, and architecture cannot serve as a
competing explanation for any transport effect observed.

Targets are five pathologies --- Cardiomegaly, Edema, Pleural Effusion,
Atelectasis, Consolidation --- labelled by a single labeller family applied
identically at both sites, so a site effect cannot be confounded with a labeller
artefact. Uncertain and unmentioned labels carry no supervision and are excluded
from the loss rather than mapped to negative; mapping them would alter the
prevalence the protocol froze.

Per-pathology thresholds are selected by balanced accuracy on the source
calibration tier. Balanced accuracy rather than accuracy because these
pathologies are far from balanced, and at Consolidation's prevalence a
plain-accuracy threshold drifts toward predicting the majority class
everywhere. Confidence is $\lvert 2p-1 \rvert$ per pathology; study confidence
is the minimum across the five, so a study is only as trustworthy as its least
certain finding. Loss is five-label Hamming error and abstention means referral
for human interpretation. Coverage is set at \SI{80}{\percent} (primary), with
\SI{90}{\percent} and \SI{70}{\percent} as sensitivity analyses, and transferred
frozen.

Coverage is held fixed and error allowed to vary, rather than the reverse, so
the two sites are compared at the same operating burden on the human reader.

\subsection{Statistical analysis}
\label{sec:power}

Inference is by patient-level bootstrap, \num{2000} replicates, seed
\num{20260718}, \SI{95}{\percent} percentile intervals. Patients are the
resampling unit because a patient contributes several studies whose errors are
correlated; resampling studies would treat those as independent and understate
the interval.

Controlled context comparisons are paired within site: C0, C1 and C2 are the
same patients under three interventions, so each replicate draws a patient set
once and reads all three conditions for it. Cross-institution comparisons use
independent within-site bootstraps, since the sites hold different patients.
Per-pathology secondary analyses use Holm correction; the primary estimand is
the aggregate Hamming loss and pays no multiplicity penalty.

A hypothesis is declared confirmed when its interval excludes zero \emph{and}
the point estimate reaches a materiality threshold of \num{0.02} absolute.

\paragraph{A note on that rule.} Because it requires the point estimate itself
to reach materiality, evaluating it at a true effect exactly equal to
materiality yields rejection in only half of repetitions: power is capped at
\SI{50}{\percent} regardless of sample size. We verified this holds at an
implied $n$ of order $10^{9}$. Materiality is therefore the smallest effect
worth acting on, not the design effect, and the study is powered against the
minimum detectable effect instead. This ceiling generalises to any study using a
``interval excludes zero \textsc{and} estimate $\geq \delta$'' decision rule.

Design-stage power was computed over a declared grid of base risks, effect
sizes, intra-patient correlations and coverages --- swept, never fitted, since
no outcome existed. On the prespecified partition the minimum detectable effect
at \SI{80}{\percent} power is \numrange{0.022}{0.026} absolute, against
\numrange{0.026}{0.072} on the official test split. The normal approximation
used for the grid agreed with the frozen bootstrap procedure to within
\num{0.015} across validation cells. Any hypothesis whose minimum detectable
effect exceeds the achievable range is reported as estimation with an interval
rather than as a powered decision.

A coverage drift of \num{0.05} from target is prespecified as material. A frozen
cutoff is intended to hold the operating burden constant across sites and
conditions; where it does not, the comparison is no longer at equal coverage,
and the drift is itself a result rather than a nuisance.

\subsection{External site}

The external site publishes permitted context as structured fields rather than
as sections parsed from a report body. The mechanism differs, the intent does
not, and the informativeness rule is the source-site rule applied without
refitting --- refitting it on external text would constitute external-site
tuning, which the registry prohibits.

Normalisation, channel handling, study aggregation and label masking are
identical to the source site. Both sites reduce full-resolution images to the
model input in a single bilinear step; applying a second resampling only at the
external site would introduce a preprocessing difference tracking the site
variable under study, and was avoided at the cost of a larger transfer.

Unlike the source site, where every cohort study is informative by construction,
the natural-state distribution is informative at the external site and is
reported.
